
Despite 160,800 human judgments showing substantial consistency in identifying safety violations (Cohen's κ=0.724), standard language model embeddings reveal no geometric structure separating safe from unsafe responses (Cohen's d=0.034, p=0.797). This disconnect between human-detectable and embedding-detectable structure has implications for alignment evaluation in RLHF-trained systems.

Reinforcement Learning from Human Feedback (RLHF) has become a paradigm for aligning large language models with human values [Christiano et al., 2017; Ouyang et al., 2022; Bai et al., 2022]. Current methods train reward models on human preference data—treating chosen versus rejected response pairs as supervision for a scalar reward function—then use these models to guide policy optimization. Once trained, the preference data itself is typically discarded.

This raises a question: does RLHF preference data encode discoverable geometric structure that could serve as reusable alignment benchmarks? If human annotators consistently identify safety violations across 160,000+ judgments, one might hypothesize that this aggregated signal induces structure in semantic embedding space—with rejected responses clustering in interpretable patterns, distances encoding violation severity, and principal component directions encoding violation types.

Our systematic investigation reveals a limitation: while humans detect alignment violations with substantial consistency (inter-annotator agreement κ=0.724), pretrained semantic encoders (RoBERTa-base) fail to capture this structure. Analyzing 160,800 chosen/rejected pairs from Anthropic's HH-RLHF dataset, we find rejected responses show no meaningful clustering (MANOVA effect size d=0.034), only 8.5× above random baseline (d=0.004). With statistical power exceeding 0.99 to detect medium effects (d≥0.5), this represents a genuine null result rather than a Type II error from insufficient data.

The key finding is that human-detectable alignment structure does not imply embedding-space structure when using standard pretrained encoders. Pretrained encoders optimize for general semantic similarity through masked language modeling objectives rather than safety-specific feature discrimination. Alignment violations are semantically diverse—spanning toxicity, misinformation, and instruction failures—and occupy overlapping embedding regions despite being consistently distinguishable to human annotators using explicit safety criteria.

This negative finding provides value by establishing that embedding-space clustering approaches using standard pretrained models are insufficient for alignment evaluation. Our validation protocol makes methodological contributions: we demonstrate that HH-RLHF contains genuine violations at 45.6\% base-rate (binomial p=0.0063), provide a framework for auditing preference dataset quality, and demonstrate how systematic hypothesis decomposition can pinpoint failure points in complex causal chains.

\textbf{Contributions:}

\begin{enumerate}
\item \textbf{Empirical negative result:} First systematic test of the geometric manifold hypothesis for RLHF alignment failures, demonstrating that standard pretrained encoders are insufficient for capturing alignment structure despite human detection consistency (κ=0.724).
\end{enumerate}

\begin{enumerate}
\item \textbf{Methodological contribution:} Base-rate validation protocol for RLHF datasets, establishing that HH-RLHF harmless subset contains 45.6\% genuine safety violations (95\% CI: [41.3\%, 50.0\%]).
\end{enumerate}

\begin{enumerate}
\item \textbf{Empirical demonstration:} Documentation of human-embedding space disconnect for alignment tasks—consistent human judgment (κ=0.724) coexists with random-like embedding distribution (d=0.034).
\end{enumerate}

\begin{enumerate}
\item \textbf{Research direction:} Evidence that geometric approaches require safety-specialized representations, with proposed directions including reward model embeddings and safety-fine-tuned encoders.
\end{enumerate}

The remainder of this paper is organized as follows: Section 2 positions this work within RLHF literature and embedding-based evaluation methods. Section 3 describes the three-hypothesis validation protocol. Section 4 presents experimental setup and evaluation criteria. Section 5 reports results for each hypothesis. Section 6 discusses implications, limitations, and future directions. Section 7 concludes.
